% ============================================================================
% Chapter 15: Basic Troubleshooting
% ============================================================================

\chapter{Basic Troubleshooting}
\label{ch:troubleshoot}

\begin{objectives}
  \item Identify common computer problems
  \item Follow a logical troubleshooting process
  \item Fix simple issues like frozen screens and no sound
  \item Know when to ask for help from a teacher or IT support
\end{objectives}

\bytetip[tip]{Things go wrong sometimes --- that's perfectly normal! The important thing is knowing how to fix them calmly and logically. Let's learn the troubleshooting mindset!}

\section{The Troubleshooting Flowchart}
\label{sec:troubleshoot-flow}
\index{troubleshooting}

When something goes wrong, follow this step-by-step process:

% --- Figure 15.1: Troubleshooting Flowchart ---
\begin{figure}[H]
\centering
\begin{tikzpicture}[node distance=0.9cm]
  \node[terminal] (s1) {Problem Occurs};
  \node[process, below=of s1] (s2) {Try Restarting\\the App};
  \node[decision, below=of s2, minimum width=2.5cm] (d1) {Fixed?};
  \node[process, below=1cm of d1] (s3) {Restart the\\Computer};
  \node[decision, below=of s3, minimum width=2.5cm] (d2) {Fixed?};
  \node[process, below=1cm of d2] (s4) {Check All Cable\\Connections};
  \node[decision, below=of s4, minimum width=2.5cm] (d3) {Fixed?};
  \node[terminal, below=1cm of d3, fill=FunSky!10, draw=FunSky] (s5) 
    {Ask Your Teacher\\/ IT Support};
  
  % Right side "Yes" exits
  \node[terminal, right=2.5cm of d1, fill=LabGreen!10, draw=LabGreen] (done1) {Done!};
  \node[terminal, right=2.5cm of d2, fill=LabGreen!10, draw=LabGreen] (done2) {Done!};
  \node[terminal, right=2.5cm of d3, fill=LabGreen!10, draw=LabGreen] (done3) {Done!};
  
  % Main flow
  \draw[thickarrow] (s1) -- (s2);
  \draw[thickarrow] (s2) -- (d1);
  \draw[labelarrow=WarnRed, line width=1.5pt] (d1.south) -- (s3) 
    node[midway, left, font=\scriptsize\sffamily\bfseries, text=WarnRed] {NO};
  \draw[thickarrow] (s3) -- (d2);
  \draw[labelarrow=WarnRed, line width=1.5pt] (d2.south) -- (s4) 
    node[midway, left, font=\scriptsize\sffamily\bfseries, text=WarnRed] {NO};
  \draw[thickarrow] (s4) -- (d3);
  \draw[labelarrow=WarnRed, line width=1.5pt] (d3.south) -- (s5) 
    node[midway, left, font=\scriptsize\sffamily\bfseries, text=WarnRed] {NO};
  
  % Yes branches
  \draw[labelarrow=LabGreen, line width=1.5pt] (d1.east) -- (done1)
    node[midway, above, font=\scriptsize\sffamily\bfseries, text=LabGreen] {YES};
  \draw[labelarrow=LabGreen, line width=1.5pt] (d2.east) -- (done2)
    node[midway, above, font=\scriptsize\sffamily\bfseries, text=LabGreen] {YES};
  \draw[labelarrow=LabGreen, line width=1.5pt] (d3.east) -- (done3)
    node[midway, above, font=\scriptsize\sffamily\bfseries, text=LabGreen] {YES};
\end{tikzpicture}
\caption{The troubleshooting flowchart --- follow each step before escalating}
\label{fig:troubleshoot}
\end{figure}

\section{Common Problems and Fixes}
\label{sec:common-problems}

\begin{table}[H]
\centering
\caption{Common computer problems and quick fixes}
\label{tab:common-fixes}
\renewcommand{\arraystretch}{1.3}
\begin{tabularx}{\textwidth}{l X X}
\toprule
\rowcolor{HeaderRow}
\textcolor{white}{\textbf{Problem}} & \textcolor{white}{\textbf{Possible Cause}} & \textcolor{white}{\textbf{Quick Fix}} \\
\midrule
\rowcolor{RowLight} Screen is black & Monitor off or cable loose & Check power \& cables \\
\rowcolor{RowDark} No sound & Volume muted or low & Click speaker icon, unmute \\
\rowcolor{RowLight} App frozen & Program crashed & \key{Ctrl}+\key{Alt}+\key{Del} $\rightarrow$ Task Manager \\
\rowcolor{RowDark} No internet & Wi-Fi disconnected & Reconnect to Wi-Fi \\
\rowcolor{RowLight} Slow computer & Too many apps open & Close unused programs \\
\bottomrule
\end{tabularx}
\end{table}

\bytetip[warn]{If your screen is completely frozen and nothing responds, hold the \key{Power} button for 5--10 seconds to force restart. But only do this as a \textbf{last resort} --- it can cause data loss!}

\begin{shortcutbox}
\begin{tabular}{ll}
  \key{Ctrl}+\key{Alt}+\key{Del} & Open security options / Task Manager \\
  \key{Ctrl}+\key{Shift}+\key{Esc} & Open Task Manager directly \\
  \key{Alt}+\key{F4} & Close the current window/program \\
\end{tabular}
\end{shortcutbox}

\begin{funfact}
The \key{Ctrl}+\key{Alt}+\key{Del} shortcut was invented by \textbf{David Bradley}, an IBM engineer, in 1981. He originally designed it as a quick way to restart during development. He once joked: ``I may have invented it, but Bill Gates made it famous!''
\end{funfact}

\begin{reviewqs}
  \item What is the first thing you should try when an app stops responding?
  \item How do you open Task Manager?
  \item Name three common computer problems and their fixes.
  \item When should you ask your teacher for help instead of trying to fix it yourself?
  \item What does \key{Ctrl}+\key{Alt}+\key{Del} do?
\end{reviewqs}

\begin{challenge}
Open \textbf{Task Manager} (\key{Ctrl}+\key{Shift}+\key{Esc}). Look at the \textbf{Processes} tab. Which program is using the most memory? Which is using the most CPU? Write down your findings!
\end{challenge}
